\documentclass[../../main.tex]{subfiles}
\begin{document}

{\bf [解]}
$C: xy=-2$，$c=\cos\theta,\ s=\sin\theta$とする．

\paragraph{(1)}
$Q(X,Y)$とおくと，これは$P$を$\theta$だけ回転した座標だから複素数の回転の性質から
\begin{align}
X+Yi=\left(t-\dfrac{2}{t}i\right)(c+is)=\left(tc+\dfrac{2}{t}s\right)+i\left(ts-\dfrac{2}{t}c\right)
\end{align}
である．$t,s,c$は実数だから実部・虚部を比較して，
\begin{align}
Q\left(t\cos\theta+\dfrac{2}{t}\sin\theta,\ t\sin\theta-\dfrac{2}{t}\cos\theta\right)
\end{align}
が求める座標である．

\paragraph{(2)}
$Q(X,Y)$を時計回りに$\theta$だけ戻すと$P$になるので，$P(x,y)$として
\begin{align}
x+yi=(X+Yi)(c-is)=(Xc+Ys)+i(-Xs+Yc)
\end{align}
だから，(1)と同様に
\begin{align}
x=Xc+Ys, \qquad y=-Xs+Yc
\end{align}
だから$C$の方程式に代入して
\begin{align}
(Xc+Ys)(-Xs+Yc)=-2
\end{align}
よって$Q$が動く曲線は
\begin{align}
 (x\cos\theta+y\sin\theta)(-x\sin\theta+y\cos\theta)=-2 \label{eq:1}
\end{align}
である．

\paragraph{(3)}

\cref{eq:1}に$(x,y)=(\sqrt{3}+1,\sqrt{3}-1)$を代入して整理すると
\begin{align}
\begin{split}
\left\{(\sqrt{3}+1)c+(\sqrt{3}-1)s\right\}\left\{-(\sqrt{3}+1)s+(\sqrt{3}-1)c\right\}&=-2 \\
-(\sqrt{3}+1)(\sqrt{3}-1)s^2+(\sqrt{3}+1)(\sqrt{3}-1)c^2+\left\{(\sqrt{3}-1)^2-(\sqrt{3}+1)^2\right\}sc&=-2 \\
2(c^2-s^2)-4\sqrt{3}cs&=-2 \\
-\cos2\theta+\sqrt{3}\sin2\theta&=+1 \\
2\sin\left(2\theta-\dfrac{\pi}{6}\right)&=1 \quad\cdots\text{①}
\end{split} \label{eq:2}
\end{align}
を得る．

$0<\theta<2\pi$より$-\dfrac{\pi}{6}<2\theta-\dfrac{\pi}{6}<\dfrac{23}{6}\pi$だから，\cref{eq:2}をみたすのは，
\begin{align}
\begin{split}
2\theta-\dfrac{\pi}{6}&=\dfrac{\pi}{6},\dfrac{5}{6}\pi,\dfrac{13}{6}\pi,\dfrac{17}{6}\pi \\
\theta&=\dfrac{1}{6}\pi,\dfrac{3}{6}\pi,\dfrac{7}{6}\pi,\dfrac{9}{6}\pi \\
&=\dfrac{1}{6}\pi,\dfrac{1}{2}\pi,\dfrac{7}{6}\pi,\dfrac{3}{2}\pi
\end{split}
\end{align}
の時である．

\newpage

\end{document}
